\documentclass[10pt,a4paper]{article}
\usepackage[T1]{fontenc}
\usepackage{lmodern,amsmath,amssymb,amsthm,microtype}
\usepackage[margin=25mm]{geometry}
\usepackage[colorlinks=true,linkcolor=blue,citecolor=blue,urlcolor=blue]{hyperref}
\setlength{\emergencystretch}{2em}
\newtheorem{proposition}{Proposition}
\newcommand{\Q}{\mathbb Q}
\title{Aligned contact numerators and physical detector scales}
\author{}\date{17 November 2026}
\begin{document}\maketitle

\begin{abstract}
An explicit slope relation makes the two general contact numerators
identical. Both the original and normalized weight ratios are then
rational squares. The physical contact normalizers prevent this
identity from identifying even the order-one detector coefficients
on the prepared positive domain. The result isolates a usable
square-class identity and the additional conditions that a shared
pole-witness construction would still have to meet.
\end{abstract}

\section*{The physical weights}
Retain the contact set of the general source \cite{Reweighted}:
\[
 \Theta=\{0,-1/2,-1,1/2399,-32/407\},\qquad
 \Theta_0=(2+\Theta)\cup(4+\Theta),
\]
and put
\[
 P(t)=\prod_{\theta\in\Theta_0}(t-\theta).
\]
Its ten roots are distinct and positive. Let $\rho,T\in\Q^\times$
with $\rho>0$, and define
\[
 a=1+\rho T^2,\qquad
 R_0=\frac{P(\lambda)}{(\lambda+1)^{10}},\qquad
 R_1=\frac{P(a\lambda-4)}{(a\lambda+2)^{10}},
\]
\[
 u_j=\rho R_j,\qquad
 w_j=\frac{u_j}{(1+u_j)^2},\qquad
 \delta_0=\lambda+1,\quad \delta_1=a\lambda+2.
\]
The scale $a$ is fixed by the actual root traces; it is not an
independent variable. The detectors contain factors of the form
$1+C^2/(\delta_j^2u_j)$, or their normalized counterparts with
$w_j$. The even-abscissa variant \cite{Even} replaces certain
numerators $C^2$ by $C$ while retaining these same denominators.

\begin{proposition}\label{prop:alignment}
Impose
\[
                   \lambda=\frac4{\rho T^2},
\]
and suppose that $\lambda$ exceeds every root of $P$ and
$0<\rho<2/3$. Put $s=(\lambda+1)/(\lambda+6)$. Then $0<s<1$,
all the displayed weights and normalizers are positive, and
\begin{align}
 a\lambda&=\lambda+4,& R_1&=s^{10}R_0,& u_1&=s^{10}u_0,\label{eq:raw}\\
 \frac{w_1}{w_0}
 &=\left(\frac{s^5(1+u_0)}{1+s^{10}u_0}\right)^2.\label{eq:norm}
\end{align}
Nevertheless, the two corresponding detector denominators satisfy
\begin{align}
 \frac{\delta_1^2u_1}{\delta_0^2u_0}&=s^8<1,\label{eq:rawscale}\\
 \frac{\delta_1^2w_1}{\delta_0^2w_0}
 &=\left(\frac{s^4(1+u_0)}{1+s^{10}u_0}\right)^2<1.\label{eq:normscale}
\end{align}
Thus this alignment does not identify the order-one detector
coefficients when the field tests are kept unchanged.
\end{proposition}
\begin{proof}
The slope relation gives $\rho T^2\lambda=4$, hence
$a\lambda=\lambda+4$. The two contact numerators are consequently
both $P(\lambda)$, while $\delta_1=\lambda+6$.
This proves \eqref{eq:raw}; substitution into the normalized
weights gives \eqref{eq:norm}. Dividing by the contact normalizers
proves the identities in \eqref{eq:rawscale}--\eqref{eq:normscale}.

For each $\theta\in\Theta_0$ one has
$0<(\lambda-\theta)/(\lambda+1)<1$. Thus
$0<R_0<1$ and $0<u_0<\rho<2/3$. Equality of the positive square
root in \eqref{eq:normscale} to one would require
\[
 u_0=u_*:=\frac{1-s^4}{s^4-s^{10}}
          =\frac{1+s^2}{s^4(1+s^2+s^4)}.
\]
The exact factorization
\[
 u_*-\frac23
 =\frac{(1-s^2)(2s^6+4s^4+6s^2+3)}
        {3s^4(1+s^2+s^4)}
\]
shows $u_*>2/3$. Moreover,
\[
 1+s^{10}u_0-s^4(1+u_0)
                 =(s^4-s^{10})(u_*-u_0)>0.
\]
This proves the strict inequality in \eqref{eq:normscale}.
The one in \eqref{eq:rawscale} follows immediately from $0<s<1$.
\end{proof}

\section*{Square classes and the remaining arithmetic conditions}
Whenever $V_0,V_1\in\Q^\times$ are the actual reciprocal traces,
the two normalized pole inputs $b_j=4w_j/V_j^2$ have the same
rational square class on this alignment. More precisely,
\[
 \frac{b_1}{b_0}
 =\left(\frac{s^5(1+u_0)V_0}
                {(1+s^{10}u_0)V_1}\right)^2.
\]
This statement uses the actual traces but imposes no relation
between them. It does not show that $(1-7b_1)/(1-7b_0)$ is a
square, that the two intrinsic pole blocks can share their full
witness pair, or that either trace-conic radicand is a square.

The coefficient inequality concerns the indicated terms with their
unchanged field tests. It does not rule out accidental equality
of the complete products at a special tuple, a more general
correspondence between pole surfaces, or individual rational
solutions on either actual square cover.

The slope is fixed once $\rho$ and the root traces have fixed $T$.
In the inherited completeness proof the roots and scale are fixed
before the fresh comparison primes and slope are chosen. Its
prime-pattern theorem therefore does not itself produce an aligned
family: that would be an additional arithmetic realization assertion.
The general scalar bounds in \cite{Scalar} give
$0<\rho<1/56$, so the small-scalar hypothesis of
Proposition~\ref{prop:alignment} is satisfied by those scalars; the
separate condition that this particular slope lie beyond the contacts
still has to be checked.

The computations in
\nolinkurl{computation/pole_parameters/general21_alignment_v75}
verify the rational identities and threshold factorization exactly.
An independent integer-polynomial convolution checks the decisive
positive factorization. No actual elliptic witness or complete
twenty-one-variable definition is asserted.

\begin{thebibliography}{9}\raggedright
\bibitem{Reweighted}\emph{Quadratic reweightings of native elliptic detectors},
\nolinkurl{papers/pole_parameters/reweighted_elliptic_detectors_20261116_v53.tex}.
\bibitem{Even}\emph{Even abscissa valuations in the general twenty-five-variable source},
\nolinkurl{papers/few_variables25/even_valuation25_20261117_v75.tex}.
\bibitem{Scalar}\emph{Intrinsic congruent-curve poles from a fixed scalar rescaling},
\nolinkurl{papers/pole_parameters/general_intrinsic_transfer_20261117_v57.tex}.
\end{thebibliography}
\end{document}
